\subsection{Intro}
\subsubsection{Groups and Group Representations}

\begin{definition}
    A group $G$ is a set with a binary operation s.t.

    $(a,b) \rightarrow a/cdot b$
    \begin{itemize}
        \item Associativity: $(a\cdot b)\cdot c = a(b \cdot c)$
        \item Identity: $\exists e \in G$ s.t. $a\cdot e = a = e \cdot a$
        \item Inverse: Given $a \in G, \; \exists b\in G s.t. ab=e=ba$
    \end{itemize}
\end{definition}

\begin{example}
    \begin{itemize}
        \item Show that the identity is unique
        \item If $ab=ac$, then $b=c$
    \end{itemize}
\end{example}

\begin{itemize}
    \item \begin{proof}
        Let $e_1,e_2$ be identities in $G$

        $e_1 = e_1e_2 = e_1$
    \end{proof}
    \item \begin{proof}
        $ab=ac$

        $\exists a^{-1} s.t. a^{-1}a=e$

        $a^{-1}(ab)=a^{-1}(ac) = b = c$
    \end{proof}
\end{itemize}

\begin{example}
    \begin{itemize}
        \item $\{e\}$
        \item $\left( \Z, + \right)$
        \item $\left( \Q^{\times}, \cdot \right)$
        \item $\Z \setminus n\Z$ (Cyclic Groups)
        
        $n \in \Z$ if $a,b \in \Z$ are congruent with $n$ if $n \mid (a-b)$

        \begin{definition}
            Write $[a]$ as the equivalence class of $a$.
            
            $[a]+[b]=[a+b]$
        \end{definition}
        Check if $[a]=[a']$ and $[b]=[b'] \implies [a+b]=[a'+b']$

        Abelian: commutative

        \item $S_3$ Permutations of $\{1,2,3\}$
        
        \begin{itemize}
            \item (1) = $e$
            \item (12) = $\sigma$
            \item (13) = $\tau\sigma=\sigma\tau^2$
            \item (23) = $\sigma\tau$
            \item (123) = $\tau$
            \item (132) = $\tau^2$
        \end{itemize}
        \begin{center}
            \begin{tabular}{ c|cccccc } 
            top first& $e$ & $\sigma$ & $\sigma\tau$ & $\sigma\tau^2$ & $\tau$ & $\tau^2$ \\
             \hline 
             $e$ & $e$ & $\sigma$ & $\sigma\tau$ & $\sigma\tau^2$ & $\tau$ & $\tau^2$ \\ 
             $\sigma$ & $\sigma$ & $e$ & $\tau$ & $\tau^2$ & $\sigma\tau$ & $\sigma\tau^2$ \\
             $\sigma\tau$ & $\sigma\tau$ & $\tau^2$ & $e$ & $\tau$ & $\sigma\tau^2$ & $\sigma$ \\
             $\sigma\tau^2$ & $\sigma\tau^2$ & $\tau$ & $\tau^2$ & $e$ & $\sigma$ & $\sigma\tau$ \\
             $\tau$ & $\tau$ & $\sigma\tau^2$ & $\sigma$ & $\sigma\tau$ & $\tau^2$ & $e$ \\
             $\tau^2$ & $\tau^2$ & $\sigma\tau$ & $\sigma\tau^2$ & $\sigma$ & $e$ & $\tau$
            \end{tabular}
            \end{center}
        This is the same as $D_3$

        \[
        e = \begin{bmatrix}
            1 & 0\\
            0 & 1
            \end{bmatrix}\\
        \]\[
        \sigma = \begin{bmatrix}
            -1 & 0\\
            0 & 1
            \end{bmatrix}\\
        \]\[
        \tau = \begin{bmatrix}
            \cos(\theta) & -\sin(\theta)\\
            \sin(\theta) & \cos(\theta)
            \end{bmatrix}
        \theta=120
        \]\[
        \tau^2 = \begin{bmatrix}
            \cos(\theta) & -\sin(\theta)\\
            \sin(\theta) & \cos(\theta)
            \end{bmatrix}
        \theta=240
        \]\[
        \sigma\tau = \begin{bmatrix}
            -\cos(\theta) & \sin(\theta)\\
            \sin(\theta) & \cos(\theta)
            \end{bmatrix}
        \theta=120
        \]\[
        \sigma\tau = \begin{bmatrix}
            -\cos(\theta) & \sin(\theta)\\
            \sin(\theta) & \cos(\theta)
            \end{bmatrix}
        \theta=240
        \]
    \end{itemize}
\end{example}
